\documentclass{article}
\usepackage{tensors}
\usepackage{braket}
\usepackage{parskip,fullpage}
\usepackage{relsize}

%% Hyperlinks
\usepackage{hyperref}
\hypersetup{
  naturalnames=false,
  hypertexnames=false,
  breaklinks,
  colorlinks = true,
  citecolor = green!50!black,
  filecolor = cyan!50!black,
  urlcolor = magenta!50!black,
  %allcolors = blue!50!black,
}

%% For smart cross-references
\usepackage[capitalize,nameinlink]{cleveref}

% For LaTeX code listings and much more
\usepackage{lmodern}
\usepackage{listings}
\usepackage[most]{tcolorbox}

% Colors for "keywords"
\colorlet{kw2}{blue} 
%\colorlet{kw3}{green!50!black} % 
%\colorlet{kw4}{orange!75!black} % 


%% -----
%% Hack to create a new "add to literate" option for listings
%% -----

%https://tex.stackexchange.com/questions/149056/how-can-i-define-additional-literate-replacements-without-deleting-existing-ones
\makeatletter
\def\addToLiterate#1{\edef\lst@literate{\unexpanded\expandafter{\lst@literate}\unexpanded{#1}}}%
\lst@Key{add to literate}{}{\addToLiterate{#1}}
\makeatother


\lstdefinestyle{tensorlatex}{
  style=tcblatex,
  commentstyle=\color{red!70!black},
  basicstyle=\ttfamily,
  texcsstyle=*,
  keywordstyle={[2]\color{kw2}},
  morekeywords={[2] 
    Vc, Mx, Tr, Tn, Tm, Inv,
    prn, angb, crly, sqr, ceil, floor, dsqr, nrm, tnrm, fnrm, abs,
    top, krn, krp, had, ttm, bigtop, bigkrn, bigkrp, bighad,bigttm,
    miwc, siwc, minc, sinc, msiz, ssiz, Rmsiz, Rssiz, mdom, sdom, msum, ssum, %dsum, dsub,
    mtop, smtop, sktop, mttm, sttm,
    mkrn, skrn, mbigkrn, sbigkrn,
    mkrp, skrp, mbigkrp, sbigkrp,
    mhad, shad, shadgram, sbighad,
  },
  % keywordstyle={[3]\color{kw3}},
  % morekeywords={[3]},
  % keywordstyle={[4]\color{kw4}},
  % morekeywords={[4]},
  add to literate={*}{\textcolor{kw2}{*}}{1},
  add to literate={+}{\textcolor{kw2}{+}}{1},
  add to literate={!}{\textcolor{kw2}{!}}{1},
}

\lstdefinestyle{tensorlatexplus}{
  style=tensorlatex,
  moredelim=**[is][\slshape]{@}{@},
  moredelim=**[is][\color{gray}]{?}{?},  
}

\lstset{
  style=tensorlatexplus,
}

% --- Code Box ---
% For reasons that I don't understand, some latexboxes must be
% followed by explicit {} to keep it from getting confused.
\DeclareTCBListing{latexbox}{s s G{1.5in} O{}}{%
  IfBooleanTF={#1}%
  {IfBooleanTF={#2}{listing side text}{text side listing}, lefthand width=#3}%
  {},%
  listing options={style=tensorlatex},%
  #4
}

\begingroup
\makeatletter
% Let's use ^^A instead of %.
% Make % other so it can be "fed" to \scantokens to handle \scantokens'
% insertion of \endlinechar;
\catcode`\^^A=14\relax
\catcode`\%=12\relax
\@firstofone{^^A
  \endgroup
  \NewDocumentCommand\ExampleRow{svmO{0pt}}{^^A
    {{\IfBooleanF{#1}{$\displaystyle}\scantokens{#2%}\IfBooleanF{#1}{$}}}&^^A
    {{\lstinline{#2}}}&^^A
    {#3}\\[#4]\hline^^A
  }^^A
}%

\title{The \texttt{tensors} package for \LaTeX2e}
\author{Tamara G. Kolda\footnote{Email: tammy.kolda@mathsci.ai}}
\date{November 10, 2024}

\begin{document}

\maketitle

\begin{abstract}
  This package provides macros for work with tensors.
  It includes special commands for formatting tensors (bold Euler Script uppercase letters),
  matrices (boldface uppercase letters), and vectors (using boldface lowercase letters) along with
  decorations, subscripts, and transpose or inverse indicators (as appropriate).
  It defines special operators for tensor outer product, Kronecker product, Khatri-Rao product (KRP),
  Hadamard product, and tensor-times-matrix (TTM) which are all equally sized and scale appropriately.
  It provides smarter containers (parentheses, angle braces, curly braces, square braces,
  ceiling braces, floor braces, double square braces, norm, and absolute value) that are
  easily resized.
  It has special operators for efficiently writing subscripts multi indices (with or without commas),
  multi sizes, multi domains, multi sums, multi tensor outer product, multi TTMs, multi Kronecker
  products, multi KRPs, and multi Hadamard products.
\end{abstract}

\section{Using this package}
\label{sec:using-this-package}

Add \lstinline|\usepackage{tensors}| to your \LaTeX\@ document header.
This package includes \texttt{xparse},
\texttt{amssymb}, \texttt{amsfonts}, \texttt{amsmath}, \texttt{amxbsy},
\texttt{eucal} (with option \texttt{mathscr}),
\texttt{stmaryrd},
\texttt{bm},
\texttt{scalerel},
\texttt{mathtools}.
It should be loaded before \lstinline{cleveref}.

% \section{General shortcut}
% \label{sec:general-shortcuts}

% Though this is not specific to tensors, we define a replacement for
% the \texttt{amsmath} command \lstinline{\text} that adds some space
% (\lstinline{\quad} by default) on either side.

% \begin{itemize}
% \item \lstinline|\qtext{@Text@}| --- Equivalent to \lstinline|\quad\text{@Text@}\quad|
% \item \lstinline|\qtext[@Space@]{@Text@}| - Replaces \lstinline{\quad} with \lstinline{@Space@} on either side of \lstinline{@Text@}
% \item \lstinline|\qtext[@LSpace@][@RSpace@]{@Text@}| - Use \lstinline{@LSpace@} to left of \lstinline{@Text@} and \lstinline{@RSpace@} to right
% \end{itemize}


\section{Tensor, Matrix, and Vector Macros}
\label{sec:vc-mx-tn}

\subsection{Tensors}
\label{sec:tensors}

The tensor command \lstinline{\Tn} makes boldface Euler script letters.
An optional argument can be used to add decoration.
A second argument in curly braces adds a subscript,
but this only works if there is no space between the first and second pair of curly braces.

\renewcommand{\arraystretch}{1.5}
\begin{tabular}{|l|l|l|} \hline
  \ExampleRow{\Tn{X}}{Tensor}
  \ExampleRow{\Tn[\bar]{X}}{Decoration}
  \ExampleRow{\Tn{X}{i}}{With subscript}
  \ExampleRow{\Tn[\tilde]{X}{i}}{With decoration and subscript}
\end{tabular}

\subsection{Matrix}
\label{sec:matrix}

The matrix command \lstinline{\Mx} is identical to the tensor command except that it
creates boldface uppercase roman letters.
The matrix command also works with Greek letters, making them bold.
A tranpose is added by following it with an apostrophe (\texttt{'}),
and a phantom transpose is produced by a quote (\texttt{"}).
The phantom transpose is useful when aligning things.

\renewcommand{\arraystretch}{1.3}
\begin{tabular}{|l|l|l|} \hline
  \ExampleRow{\Mx{A}}{Matrix}
  \ExampleRow{\Mx{A}'}{Matrix transpose}
  \ExampleRow{\Mx{\Theta}}{Greek letter}
  \ExampleRow{\Mx[\tilde]{A}}{Decoration}
  \ExampleRow{\Mx{A}{1}}{In a sequence}
  \ExampleRow{\Mx{A}{i}'\Mx{A}{i}}{Subscripts misaligned}
  \ExampleRow{\Mx{A}{i}'\Mx{A}{i}"}{Aligned using phantom tranpose}
\end{tabular}

\subsection{Matricized Tensor}
\label{sec:matricized-tensor}



The matricized tensor \lstinline|\Tm| command works exactly the same as the matrix command except that the second argument is required and
has parentheses around it.


\renewcommand{\arraystretch}{1.5}
\begin{tabular}{{|l|l|l|}}
  \hline
  \ExampleRow{\Tm{X}{k}}{Unfolded tensor in mode $k$}
  \ExampleRow{\Tm[\tilde]{X}{k}}{With decoration}
  \ExampleRow{\Tm{X}{\ell}' \text{ or } \Tm{X}{\ell}"}{With transpose or phantom transpose}
  \ExampleRow{\Tm{A}{X \times Y}}{Arbitrary subscript}
  % \ExampleRow{\Tm{A}[*]{X \times Y}}{Bigger parentheses, if you really want them}
\end{tabular}


\subsection{Vector}
\label{sec:vector}

The vector command \lstinline{\Vc} is identical to the matrix command except that it
creates boldface lowercase roman letters.

\renewcommand{\arraystretch}{1.3}
\begin{tabular}{|l|l|l|} \hline
  \ExampleRow{\Vc{v}}{Vector}
  \ExampleRow{\Vc{v}'}{Vector transpose}
  \ExampleRow{\Vc{\theta}}{Greek letter}
  \ExampleRow{\Vc[\tilde]{v}}{Decoration}
  \ExampleRow{\Vc{v}{1}}{In a sequence}
  \ExampleRow{\Vc{v}{i}'\Vc{v}{i}}{Subscripts misaligned}
  \ExampleRow{\Vc{v}{i}'\Vc{v}{i}"}{Aligned using phantom tranpose}
\end{tabular}


\subsection{Transpose \& Inverse}
The macros above for vector and matrix add a tranpose if followed by \texttt{'}.
If you need to add it manually, however, use \lstinline|^{\Tr}|.
Add \lstinline|^{\Tr*}| for a phantom transpose (useful in combination with subscripts).
For matrix inverse, use \lstinline|^{\Inv}|.
Add \lstinline|^{\Inv*}| for a phantom inverse (useful in combination with subscripts).


\renewcommand{\arraystretch}{1.3}
\begin{tabular}{|l|l|l|} \hline
  \ExampleRow{\Mx{A}^{\Tr}}{Manual transpose}
  \ExampleRow{\Mx{A}{i}^{\Tr}\Mx{A}{i}^{\Tr*}}{Aligned using phantom tranpose}
  \ExampleRow{\Mx{A}^{\Inv}}{Inverse}
  \ExampleRow{\Mx{A}{i}^{\Inv}\Mx{A}{i}^{\Inv*}}{Aligned using phantom inverse}
\end{tabular}



\section{Smarter ``Containers"}

% The \texttt{braket} package provides some smart containers
% such as \lstinline|\set| and \lstinline|\Set| (expanding version).
% \begin{latexbox}*{2in}
% $\Set{i | i \in [n]}$ ($[n] \equiv \set{1,2,\dots,n}$)
% \end{latexbox}

We define convenient shorthands for containers with the ability to resize them.

\subsection{List of Containers}
\label{sec:list-containers}



A list of containers follows below.
In some cases, we allow ``natural'' containers as well as the usual curly braces.
For example, \lstinline{\prn} is parentheses, and it can be called as \lstinline|\prn{x}|
or \lstinline{\prn(x)}.
Additionally, for the parentheses, following by an apostrophe \lstinline{'} indicates to add
transpose.

\renewcommand{\arraystretch}{1.3}
\begin{tabular}{|l|l|l|} \hline
  \ExampleRow{\prn{\cdot}, \prn(\cdot), \prn(\cdot)'}{Parentheses}
  \ExampleRow{\angb{\cdot}, \angb<\cdot>}{Angle braces} 
  \ExampleRow{\crly{\cdot}}{Curly braces}
  \ExampleRow{\sqr{\cdot},\sqr[\cdot]}{Square braces}
  \ExampleRow{\ceil{\cdot}}{Ceiling braces}
  \ExampleRow{\floor{\cdot}}{Floor braces}
  \ExampleRow{\dsqr{\cdot}}{Double square braces}
  \ExampleRow{\nrm{\cdot}}{Norm}
  \ExampleRow{\tnrm{\cdot}}{2-Norm}
  \ExampleRow{\fnrm{\cdot}}{Frobenious Norm}
  \ExampleRow{\abs{\cdot},\abs|\cdot|}{Absolute value}
\end{tabular}


\subsection{Resizing Containers}
\label{sec:resizing-containers}



One advantage of these containers is that they can be easily resized.
Adding asterisks increases the sizes using left/right versions of
\texttt{big}, \texttt{Big}, \texttt{bigg}, and \texttt{Bigg}.
Adding a plus indicates automatic sizing using \lstinline|\left| and \lstinline{\right}.
% Capitalizing the first letter switches to auto-expansion.



\renewcommand{\arraystretch}{2.5}
\begin{tabular}{|l|l|l|} \hline
  \ExampleRow{\prn{\sum_{i=1}^n x_i^2}}{Matched parentheses shorthand}[3mm]
  \ExampleRow{\prn*{\sum_{i=1}^n x_i^2}}{With left/right \texttt{big}}[3mm]
  \ExampleRow{\prn**{\sum_{i=1}^n x_i^2}}{With left/right \texttt{Big}}[3mm]
  \ExampleRow{\prn***{\sum_{i=1}^n x_i^2}}{With left/right \texttt{bigg}}[3mm]
  \ExampleRow{\prn****{\sum_{i=1}^n x_i^2}}{With left/right \texttt{Bigg}}[3mm]
  \ExampleRow{\prn+{\sum_{i=1}^n x_i^2}}{Automatic size}[3mm]
\end{tabular}


\section{Special operators}
\label{sec:matrix-operators}

We define operators for tensor outer product (\lstinline{\top}),
Kronecker product (\lstinline{\krn}), Khatri-Rao product (\lstinline{\krp}),
and tensor-times-matrix (\lstinline{\ttm}).
Here, \lstinline{\top} is \lstinline|\circ| sized to match the others operators.
The commands \lstinline{\krn} and \lstinline{\krp}
are synonyms for \lstinline|\otimes| and \lstinline|\odot|
along with their related \lstinline|big| variations that support limits.
The tensor-times-matrix \lstinline|\ttm| operation takes an optional curly-brace argument to indicate the mode; the default is $k$.

\renewcommand{\arraystretch}{1.3}
\begin{tabular}{|l|l|l|}
  \hline
  \ExampleRow{\Vc{a} \top \Vc{b}}{Tensor outer product}
  \ExampleRow{\Vc{a} \krn \Vc{b}}{Kronecker product}
  \ExampleRow{\Mx{a} \krp \Mx{b}}{Khatri-Rao product}
  \ExampleRow{\Tn{X} \ttm \Mx{U}}{Tensor-times-matrix (TTM)}
  \ExampleRow{\Tn{X} \ttm[1] \Mx{U}}{TTM with specified mode}
\end{tabular}

All of these operators have big versions that accept limits:
\lstinline{\bigtop},
\lstinline{\bigkrn},
\lstinline{\bigkrp},
\lstinline{\bigttm}.

\renewcommand{\arraystretch}{2.5}
\begin{tabular}{|l|l|l|} \hline
  \ExampleRow{\bigtop_{k=1}^d \Vc{a}{k}}{Big tensor outer product}[3mm]
  \ExampleRow{\bigkrn_{k=1}^d \Mx{A}{k}}{Big Kronecker product}[3mm]
  \ExampleRow{\bigkrp_{k=1}^d \Mx{A}{k}}{Big Khatri-Rao product}[3mm]
  \ExampleRow{\bigttm_{k=1}^d \Mx{A}{k}}{Big TTM}[3mm]
\end{tabular}


\section{Multidimensional indices, sizes, sums, and operators}
\label{sec:mult-sizes-sums}

\subsection{Multi indices with commas}
\label{sec:multi-indices-with-commas}

The \lstinline|\miwc| (multi index with commas) macro is an efficient
way to write out a $d$-way index.
The compact mode \lstinline{\miwc!} removes the second index.
The first optional argument changes the main argument from $i$ to whatever is specified.
The second optional argument changes the last index.
It's possible to do just three (one asterisk) or four (two asterisks) modes.

\renewcommand{\arraystretch}{1.3}
\begin{tabular}{|l|l|l|} \hline
  \ExampleRow{\miwc}{Multi index with comma default}
  \ExampleRow{\miwc!}{Compact}
  \ExampleRow{\miwc[j]}{Change variable}
  \ExampleRow{\miwc[j][h]}{Change last subscript}
  \ExampleRow{\miwc![j][6]}{Combination of options}
  \ExampleRow{\miwc*}{3-way}
  \ExampleRow{\miwc**[j]}{4-way}
\end{tabular}

We also define a \emph{skip} multi index using \lstinline|\siwc| which
skips a single interior index.
The missing element can be replaced using an optional curly-brace last argument.
There are three optional arguments specifying the skipped index (default $k$),
the variable (default $i$), and the last index (default $d$), respectively.

\renewcommand{\arraystretch}{1.3}
\begin{tabular}{|l|l|l|} \hline
  \ExampleRow{\siwc}{Skip index with comma default}
  \ExampleRow{\siwc[\ell]}{Change skipped index}
  \ExampleRow{\siwc[k][j]}{Change variable}
  \ExampleRow{\siwc[k][i][h]}{Change last index}
  \ExampleRow{\siwc{j_k}}{Replace skipped entry}
\end{tabular}

The commands \lstinline|\miwc| and \lstinline|\siwc| are not limited
to indices. For instance, we can use it to list out a sequence of
matrices.

\renewcommand{\arraystretch}{1.3}
\begin{tabular}{|l|l|l|} \hline
  \ExampleRow{\miwc[\Mx{A}]}{Using matrices}
  \ExampleRow{\siwc[k][\Mx{A}]{\Mx{B}_k}}{With skip}
\end{tabular}

Double subscripting can cause errors, so we create a special macro for
that \lstinline|\dsub| that works as follows that can be passed as the
first argument the \lstinline{\miwc} and similar macros.

\renewcommand{\arraystretch}{1.3}
\begin{tabular}{|l|l|l|} \hline
  \ExampleRow{\miwc[\dsub{i}{\pi}]}{Double index macro}
\end{tabular}


\subsection{Multi indices without commas}
\label{sec:multi-indic-with}


The \lstinline|\minc| (multi index no commas) is analogous to \lstinline|\miwc|
except there are no commas and the dots are centered vertically.

\begin{tabular}{|l|l|l|} \hline
  \ExampleRow{\minc}{Multi index no comma default}
  \ExampleRow{\minc!}{Compact}
  \ExampleRow{\minc[j]}{Change variable}
  \ExampleRow{\minc[j][h]}{Change last subscript}
  \ExampleRow{\minc![j][6]}{Put it together}
  \ExampleRow{\minc*}{3-way}
  \ExampleRow{\minc**[j]}{4-way}
\end{tabular}

The \lstinline|\sinc| (skip multi index no commas)
is likewise analogous to \lstinline|\siwc|.

\begin{tabular}{|l|l|l|} \hline
  \ExampleRow{\sinc}{Skip index no comma default}
  \ExampleRow{\sinc[\ell]}{Change skipped index}
  \ExampleRow{\sinc[k][j]}{Change variable}
  \ExampleRow{\sinc[k][i][h]}{Change last index}
  \ExampleRow{\sinc{j_k}}{Replace skipped entry}
\end{tabular}

\subsubsection{Multi  sizes}
\label{sec:sizes}

The \lstinline|\msiz| macros 
is analogous to \lstinline|\miwc|
except that default variable is $n$, the default separator is $\times$,
and the dots are centered.

\begin{tabular}{|l|l|l|} \hline
  \ExampleRow{\msiz}{Multi size default}
  \ExampleRow{\msiz!}{Compact}
  \ExampleRow{\msiz[m]}{Change variable}
  \ExampleRow{\msiz[n][h]}{Change last subscript}
  \ExampleRow{\msiz![m][6]}{Put it together}
  \ExampleRow{\msiz*}{3-way}
  \ExampleRow{\msiz**[m]}{4-way}
\end{tabular}

Additionally,
it's possible to provide your own letters, separated by commas, from two up to four.
If a second single letter is provided, it's seperated also by dots.

\renewcommand{\arraystretch}{1.3}
\begin{tabular}{|l|l|l|} \hline
  \ExampleRow{\msiz{m,n}}{Two arguments}
  \ExampleRow{\msiz{m,n,p}}{Three arguments}
  \ExampleRow{\msiz{m,n,p,q}}{Four arguments}
  \ExampleRow{\msiz{m}{z}}{Two groups (second group must be singular)}
  \ExampleRow{\msiz{m,n,p,q}{z}}{Two groups (first group can have up to four elements)}
\end{tabular}

The \lstinline|\ssiz| macro is likewise analogous to \lstinline|\siwc|.

\begin{tabular}{|l|l|l|} \hline
  \ExampleRow{\ssiz}{Skip size default}[1mm]
  \ExampleRow{\ssiz[\ell]}{Change skipped index}[1mm]
  \ExampleRow{\ssiz[k][m]}{Change variable}[1mm]
  \ExampleRow{\ssiz[k][n][h]}{Change last index}[1mm]
  \ExampleRow{\ssiz{m_k}}{Replace skipped entry}[1mm]
\end{tabular}


The commands \lstinline{\Rmsiz} and \lstinline{\Rssiz} add
$\mathbb{R}$ in front and make the size a subscript but are otherwise
identical to \lstinline{\msiz} and \lstinline{\ssiz}.

\renewcommand{\arraystretch}{1.5}
\begin{tabular}{|l|l|l|} \hline
  \ExampleRow{\Rmsiz}{Multi size default}
  \ExampleRow{\Rmsiz!}{Compact}
  \ExampleRow{\Rmsiz![m][6]}{With options}
  \ExampleRow{\Rmsiz*}{3-way}
  \ExampleRow{\Rmsiz**[m]}{4-way}
  \ExampleRow{\Rssiz}{Skip size default}
  \ExampleRow{\Rmsiz{n}}{One argument}
  \ExampleRow{\Rmsiz{m,n,p,q}}{Four arguments}
  \ExampleRow{\Rmsiz{n}{n}}{Two groups (second group must be singular)}
  \ExampleRow{\Rmsiz{m,n,p,q}{z}}{Two groups (first group can have up to four elements)}
\end{tabular}

\subsection{Multi  domains}
\label{sec:domains}

The \lstinline|\mdom| and \lstinline|\sdom| macros 
are essentially identical to \lstinline|\msiz| and \lstinline|\ssiz|
except that they put square brackets around the letters
and the separator is $\otimes$.

{
  \renewcommand{\arraystretch}{1.5}
\begin{tabular}{|l|l|l|} \hline
  \ExampleRow{\mdom}{Multi domain default}
  \ExampleRow{\mdom!}{Compact}
  \ExampleRow{\mdom[m]}{Change variable}
  \ExampleRow{\mdom[n][h]}{Change last subscript}
  \ExampleRow{\mdom![m][6]}{Put it together}
  \ExampleRow{\mdom*}{3-way}
  \ExampleRow{\mdom**[m]}{4-way}
  \ExampleRow{\mdom{m,n,p,q}}{Specific letters}
  \ExampleRow{\mdom{n,n}{n}}{Specific letters with skip}
  \ExampleRow{\sdom}{Skip domain default}
  \ExampleRow{\sdom[\ell]}{Change skipped index}
  \ExampleRow{\sdom[k][m]}{Change variable}
  \ExampleRow{\sdom[k][n][h]}{Change last index}
  \ExampleRow{\sdom{m_k}}{Replace skipped entry}
\end{tabular}
}

\subsection{Multi sums}
\label{sec:sums}

The \lstinline|\msum| (multi sum) and \lstinline|\ssum| macros are roughly analogous to the \lstinline|\msiz| and \lstinline|\ssiz| macros.
These provide flexible ways to write multiple sums or even single sums efficiently.
There are three optional arguments that can be used to swap: the
letter for index ($i$), the letter for the final value ($n$), and the
letter/number for the subscript of the last index ($d$).
As with \lstinline|\msiz|, we can specify only three (one asterisk) or four (two asterisks) modes.

\renewcommand{\arraystretch}{2.5}
\begin{tabular}{|l|l|l|} \hline
  \ExampleRow{\msum \mbox{ and } \msum!}{Multi sum default \& compact form}[3mm]
  %\ExampleRow{\msum!}{Compact}[3mm]
  \ExampleRow{\msum[j]}{Swap index (default is $i$)}[3mm]
  \ExampleRow{\msum[i][r]}{Swap max value (default is $n$)}[3mm]
  \ExampleRow{\msum[i][n][6]}{Swap final subscript (default is $d$)}[3mm]
  \ExampleRow{\msum* \mbox{ or } \msum**}{Three or four sums}[3mm]
  \ExampleRow{\msum*[j][r]}{Three sums with options}[3mm]
\end{tabular}


The most powerful option with \lstinline|\msiz| is the ability to
quickly specify relatively arbitary summations by providing a comma-separated list.
Each item in the list can be an index and max value pair separated by a forward slash or
just a single item that will go at the bottom of the summation.
Additionally, if a second argument is included, that indicates a separation of centered
dots before a last summation.


\renewcommand{\arraystretch}{2.5}
\begin{tabular}{|l|p{3in}|l|} \hline
  \ExampleRow{\msum{i/m,j/n,k/p}}{Arbitrary indices and max values}[3mm]
  \ExampleRow{\msum{i\in\mathcal{I}, j\in\mathcal{J}, k\in\mathcal{K}}}{Just the bottoms}[3mm]
  \ExampleRow{\msum{i_1/n,i_2/n}{i_d/n}}{Extra dots}[3mm]
\end{tabular}

\subsection{Multi tensor outer product}
\label{sec:multi-tensor-outer}

The \lstinline|\mtop| (multi tensor outer product) and \lstinline|\smtop| commands are analogous to the \lstinline|\msiz| and \lstinline|\ssiz| commands.
(Note that we use \lstinline|\sktop| rather than \lstinline|\stop| because the latter was already defined.)

\renewcommand{\arraystretch}{1.3}
\begin{tabular}{|l|l|p{1.95in}|}
  \hline
  \ExampleRow{\mtop}{Multiple tensor outer product}
  \ExampleRow{\mtop!}{Compact}
  \ExampleRow{\mtop[\Vc{b}]}{Change main letter}
  \ExampleRow{\mtop[\Vc{a}][d+1]}{Change last subscript}
  \ExampleRow{\mtop*}{Three-way}
  \ExampleRow{\mtop**}{Four-way}    
  \ExampleRow{\mtop{\Vc{a},\Vc{a}}{\Vc{a}}}{Specify first (up to 4 items) and last arguments}
  \ExampleRow{\mtop{\Vc{a},\Vc{b},\Vc{c}}}{Give specific arguments}
  \ExampleRow{\sktop}{Skip tensor outer product}
\end{tabular}



\subsection{Multi TTMs}
\label{sec:multi-ttms}

The \lstinline|\mttm| and \lstinline|\sttm| macros 
are analogous to \lstinline|\minc| and \lstinline|\sinc|
except that default variable is $\ttm \Mx{U}{k}$ where the $k$ is replaced
by the appropriate index.


\begin{tabular}{|l|l|l|} \hline
  \ExampleRow{\mttm}{Multi TTM default}
  \ExampleRow{\mttm![\Mx{V}][6]}{Some options}
  \ExampleRow{\mttm**[\Mx{V}]}{4-way}
  \ExampleRow{\sttm}{Skip TTM default}
  \ExampleRow{\sttm[\ell]{\Mx{V}}}{Some options}
\end{tabular}

\subsection{Multi Kronecker product}
\label{sec:multi-kron-prod}

We define the \lstinline|\mkrn| and \lstinline|\skrn| commands for Kronecker products.
By default, the indices go in \emph{reverse} from $d$ to 1.
They are otherwise mostly analogous to \lstinline|\miwc| and \lstinline|\siwc| in terms of options and arguments.

\begin{tabular}{|l|l|l|} \hline
  \ExampleRow{\mkrn}{Multi Kronecker default}
  \ExampleRow{\mkrn!}{Compact}
  \ExampleRow{\mkrn[\Mx{V}]}{Change main letter}
  \ExampleRow{\mkrn[\Mx{U}][d+1]}{Change last subscript}
  \ExampleRow{\mkrn*}{Three-way}
  \ExampleRow{\mkrn**}{Four-way}    
  \ExampleRow{\skrn}{Skip Kronecker default}
  \ExampleRow{\skrn{\Mx{V}}}{Replace skipped item}
  \ExampleRow{\skrn[\ell][\Mx{W}][h]}{Change options}
\end{tabular}

\subsection{Multi KRP}
\label{sec:multi-krp}

We define the \lstinline|\mkrp| and \lstinline|\skrp| commands for KRPs
which are analogous to the \lstinline|\mkrp| and \lstinline|\skrp| commands.

\renewcommand{\arraystretch}{1.5}
\begin{tabular}{|l|l|l|} \hline
  \ExampleRow{\mkrp}{Multi KRP default}
  \ExampleRow{\mkrp![\Mx{B}][d+1]}{Various options}
  \ExampleRow{\mkrp**}{Four-way}    
  \ExampleRow{\skrp}{Skip KRP default}
  \ExampleRow{\skrp{\Mx{C}}}{Replace skipped item}
  \ExampleRow{\skrp[\ell][\Mx{B}][h]}{Various options}
\end{tabular}

\subsection{Multi Hadamard}
\label{sec:multi-hadamard}

We define \lstinline{\shad} command for the multi-Hadmard product, usually of gram matrices.


{\relscale{0.85}
\renewcommand{\arraystretch}{1.5}
\begin{tabular}{|l|l|l|} \hline
  \ExampleRow{\mhad}{Multi-Hadamard}
  \ExampleRow{\mhad[\shadgram]}{Multi-Hadamard of Grams}
  \ExampleRow{\shad}{Skip multi-Hadamard}
  \ExampleRow{\shad[\ell]}{Change skip index}
\end{tabular}
}

We include the definition of a helper macro \lstinline{\shadgram} that
is a gram product of subscripted matrices.

{\relscale{0.81}
\renewcommand{\arraystretch}{1.5}
\begin{tabular}{|l|l|l|} \hline
  \ExampleRow{\shad[k][\shadgram[\Mx{B}]]}{Change letter}[0.5ex]
  \ExampleRow{\shad[k][\shadgram[\Mx{A}][\Mx{B}]]}{Two different letters!}[0.5ex]
  \ExampleRow{\shad[k][\shadgram][h]}{Change last index}[0.5ex]
\end{tabular}
}

\end{document}

%%% Local Variables:
%%% mode: latex
%%% TeX-master: t
%%% End:
